\documentclass[10pt,a4paper]{amsart}
\usepackage[margin=19mm]{geometry}
\usepackage{amsmath,amssymb,mathtools}
\usepackage[scr=rsfs,cal=euler]{mathalfa}
\usepackage{xfrac}
\usepackage[unicode,hidelinks,bookmarks=false]{hyperref}
\newcommand{\ie}{i.e.\ }
\newcommand{\cf}{cf.\ }
\newcommand{\resp}{respectively\ }
\newcommand{\Subsection}{\subsection}
\providecommand{\nobreakdash}{\nobreak\hbox{-}}
\setlength{\emergencystretch}{2em}
\multlinegap0pt
\usepackage{amssymb}
\usepackage{cancel}
\usepackage[shortlabels]{enumitem}
\newtheorem{lemma}{Lemma}
\title{The best constant in the G-N inequality for the mixed local and Nonlocal Laplacian}
\author{Hichem Hajaiej, Yu Su}
\date{Source: arXiv:2604.05177v1 (CC BY 4.0)}
\begin{document}
\maketitle

\noindent\emph{Source and licence.} The best constant in the G-N inequality for the mixed local and Nonlocal Laplacian. Version arXiv:2604.05177v1 (CC BY 4.0), distributed by arXiv under the Creative Commons Attribution 4.0 International licence, http://creativecommons.org/licenses/by/4.0/, as recorded in the arXiv licence metadata for this exact version and shown on the abstract page https://arxiv.org/abs/2604.05177v1. Full licence notice: \texttt{licenses/arxiv-2604.05177-CC-BY-4.0.txt}.\par\smallskip

\noindent\textit{Editorial scope.} Counted unit: the article's lemma Lemma4.1 (source0.tex lines 1911--1952). The article's complete original proof of it is delivered with it (source0.tex lines 3661--3868, one \texttt{proof} environment of 208 lines, which the article places away from the statement). No mathematics is written, repaired or re-derived: the statement and the proof are the article's own lines, in the article's own notation, and no retained line is substituted or re-flowed.  Retained uncounted as the source-local context the proof needs: lines 97--103; lines 739--741.  Nothing was omitted from any retained span, so the package carries no omission record.  No retained line contains a citation, so the wrapper carries no bibliography and no bibliography entry.  No retained line contains a figure environment, an image-inclusion command or drawing code, so no image is delivered and the package declares no asset dependency.  Editorial formatting only, in addition: an AMS \texttt{amsart} preamble that restates the article's own declarations and macros used by the retained lines, namely the environments \texttt{lemma} on separate counters, together with the standard packages those lines need.  The retained environments print in this document as Lemma 1 in order of appearance, whereas the article prints the same items with the numbering of its own full document.  The complete native source of the article as distributed in the arXiv e-print for this exact version is bundled unchanged as \texttt{sources/197936/source0.tex}.
\begin{equation}\label{P}
\begin{aligned}
-\Delta u
+(-\Delta)^{s}u =|u|^{p-2}u, \ \ x\in\mathbb{R}^{N},
\end{aligned}
\tag{$P$}
\end{equation}

\begin{lemma}\label{Lemma2.2}
$E\hookrightarrow L^{t}(\mathbb{R}^{N})$, $t\in [\frac{2N}{N-2s},\frac{2N}{N-2}]$.
\end{lemma}

\begin{lemma}\label{Lemma4.1}
Let $N\geqslant3$ and $0<s<1$.
Then we have

\begin{enumerate}
\item [(1)]
the functional $I$ has mountain pass geometric structure;


\item [(2)]
for any $u\in E\setminus\{0\}$, there exists a unique $t_{u}>0$ such that $t_{u}u\in \mathcal{N}$ and $I(t_{u}u)=\max\limits_{t>0}I(tu)$;

\item [(3)]
 $\bar{c}=\inf\limits_{u\in\mathcal{N}}I(u)>0$;

\item [(4)]
$c=\bar{c}=\bar{\bar{c}}$,
where
\begin{equation*}
\begin{aligned}
\bar{\bar{c}}:=\inf\limits_{u\in E\setminus\{0\}}\sup\limits_{t\geqslant 0}I(tu).
\end{aligned}
\end{equation*}

\item [(5)]
For $u\in \mathcal{N}$, we have $\Phi'(u)\not=0$, where
\begin{equation}\label{4.1}
\begin{aligned}
\Phi(u)=\langle I'(u),u\rangle=\|u\|_{E}^{2}-\int_{\mathbb{R}^{N}}|u|^{p}\mathrm{d}x,
\end{aligned}
\end{equation}
and
\begin{equation}\label{4.2}
\begin{aligned}
\langle\Phi'(u),u\rangle
=&2\|u\|_{E}^{2}
-p\beta\int_{\mathbb{R}^{N}}|u|^{p}\mathrm{d}x.
\end{aligned}
\end{equation}
Moreover, if $\bar{u}\in \mathcal{N}$ and $I(\bar{u})=c$, then $u$ is a ground state solution for equation \eqref{P}.
\end{enumerate}
\end{lemma}

\begin{proof}[Proof of Lemma \ref{Lemma4.1}]
(1) Using Lemma \ref{Lemma2.2}, one has
\begin{equation*}
\begin{aligned}
I(u)
\geqslant \|u\|_{E}^{2}
-C\|u\|_{E}^{p}.
\end{aligned}
\end{equation*}
We keep in mind that $2<p$. Then there exists a sufficiently small positive number $\rho$ such that
\begin{equation*}
\begin{aligned}
\varsigma :=
\inf_{\|u\|_{E}=\rho}
I(u)>0=I(0).
\end{aligned}
\end{equation*}
For $u\in E\setminus\{0\}$, we have
\begin{equation*}
\begin{aligned}
I(tu)
=
\frac{t^{2}}{2}
\|u\|_{E}^{2}
-\frac{t^{p}}{p}
\int_{\mathbb{R}^{N}}|u|^{p}\mathrm{d}x.
\end{aligned}
\end{equation*}
From $2<p$, it follows that $I(tu)<0$ for $t$ large enough.

From above,
we can choose $t_{1}>0$ corresponding to $u$ such that $I(t_{1}u)<0$ for $t>t_{1}$ and $\|t_{1}u\|_{E}>\rho$.

(2)
For any $u\in E\setminus\{0\}$ and $t\in (0,\infty)$, we define
\begin{equation*}
\begin{aligned}
f_{1}(t)=I(tu)=&\frac{t^{2}}{2}\|u\|_{E}^{2}-\frac{t^{p}}{p}\int_{\mathbb{R}^{N}}|u|^{p}\mathrm{d}x.
\end{aligned}
\end{equation*}
We compute
\begin{equation*}
\begin{aligned}
f_{1}'(t)
=&t\|u\|_{E}^{2}
-t^{p-1}
\int_{\mathbb{R}^{N}}|u|^{p}\mathrm{d}x.
\end{aligned}
\end{equation*}
We know that $f_{1}'(\cdot)=0$ iff
\begin{equation*}
\begin{aligned}
\|u\|_{E}^{2}
=t^{p-2}
\int_{\mathbb{R}^{N}}|u|^{p}\mathrm{d}x.
\end{aligned}
\end{equation*}
Let
\begin{equation*}
\begin{aligned}
f_{2}(t)
=t^{p-2}
\int_{\mathbb{R}^{N}}|u|^{p}\mathrm{d}x.
\end{aligned}
\end{equation*}
Clearly, $\lim\limits_{t\rightarrow0}f_{2}(t)\rightarrow0$, $\lim\limits_{t\rightarrow+\infty}f_{2}(t)\rightarrow+\infty$ .
By using the intermediate value theorem, we know that there exists a $0<t_{u}<\infty$ such that
\begin{equation*}
\begin{aligned}
f_{2}(t_{u})=\|u\|_{E}^{2}.
\end{aligned}
\end{equation*}
Furthermore, it is easy to see that $f_{2}(\cdot)$ is strictly increasing on $(0,\infty)$.
Then we get the uniqueness of $t_{u}$.
And then,
\begin{equation*}
\begin{aligned}
\|u\|_{E}^{2}
=t_{u}^{p-2}
\int_{\mathbb{R}^{N}}|u|^{p}\mathrm{d}x,
\end{aligned}
\end{equation*}
which gives
\begin{equation*}
\begin{aligned}
\|t_{u}u\|_{E}^{2}
=\beta\int_{\mathbb{R}^{N}}|t_{u}u|^{p}\mathrm{d}x
+\int_{\mathbb{R}^{N}}|t_{u}u|^{2^{*}}\mathrm{d}x.
\end{aligned}
\end{equation*}
This implies that $t_{u}u\in \mathcal{N}$.


(3)
By applying $\langle I'(u),u\rangle=0$, we know
\begin{equation*}
\begin{aligned}
0
=\langle I'(u),u\rangle\geqslant \|u\|_{E}^{2}
-C\|u\|_{E}^{p},
\end{aligned}
\end{equation*}
which implies
\begin{equation*}
\begin{aligned}
C\|u\|_{E}^{p-2}
\geqslant
1,
\end{aligned}
\end{equation*}
and
\begin{equation*}
\begin{aligned}
\|u\|_{E}^{2}
\geqslant C.
\end{aligned}
\end{equation*}
Then, for $u\in \mathcal{N}$, we get
\begin{equation*}\label{36}
\begin{aligned}
I(u)
=&
I(u)-\frac{1}{p}\langle I'(u),u\rangle\\
=&\frac{1}{2}\|u\|_{E}^{2}
-\beta\frac{1}{p}\int_{\mathbb{R}^{N}}|u|^{p}\mathrm{d}x
-\frac{1}{p}
\left(
\|u\|_{E}^{2}
-\beta\int_{\mathbb{R}^{N}}|u|^{p}\mathrm{d}x
\right)\\
=&
\left(\frac{1}{2}-\frac{1}{ p}\right)
\|u\|_{E}^{2}
\geqslant
C.
\end{aligned}
\end{equation*}
Hence, $I$ is bounded from below on $\mathcal{N}$. And then $\bar{c}>0$.





(4)
By virtue of Lemma \ref{Lemma4.1} (2), we have the following result directly
\begin{equation*}
\begin{aligned}
\bar{c}=\bar{\bar{c}}.
\end{aligned}
\end{equation*}
For any $u\in E\setminus\{0\}$,
there exists some $\tilde{t}>0$ large,
such that $I(\tilde{t}u)<0$.
Define a path $\gamma:[0, 1]\to E$ by $\gamma(t) = t\tilde{t}u$.
Clearly, $\gamma\in\Gamma$ and
\begin{equation*}
\begin{aligned}
c\leqslant \bar{\bar{c}}.
\end{aligned}
\end{equation*}
On the other hand, for each path $\gamma\in\Gamma$,
let $g(t):=\langle I'(\gamma(t)),\gamma(t)\rangle$.
Then, $g(0)=0$ and $g(t)>0$ for $t>0$ small. A direct calculation gives
\begin{equation*}
\begin{aligned}
I(\gamma(1))
-\frac{1}{p}\langle I'(\gamma(1)),\gamma(1)\rangle
\geqslant\left(\frac{1}{2}-\frac{1}{p}\right)\|\gamma(1)\|_{E}^{2}
\geqslant 0,
\end{aligned}
\end{equation*}
which implies
\begin{equation*}
\begin{aligned}
\langle I'(\gamma(1)),\gamma(1)\rangle
\leqslant p\cdot I(\gamma(1))
=p\cdot I(\tilde{t}u)<0.
\end{aligned}
\end{equation*}
Thus, there exists $\tilde{\tilde{t}} \in(0,1)$ such that $g(\tilde{\tilde{t}}) = 0$, i.e. $\gamma(\tilde{\tilde{t}})\in \mathcal{N}$ and $c\geqslant \bar{c}$.
This deduces
$c=\bar{c}=\bar{\bar{c}}$.

(5)
For $u\in \mathcal{N}$, it follows from \eqref{4.1} and \eqref{4.2} that
\begin{equation*}
\begin{aligned}
\langle\Phi'(u),u\rangle
=&\langle\Phi'(u),u\rangle-p\Phi(u)\\
=&\left(2\|u\|_{E}^{2}-p\int_{\mathbb{R}^{N}}|u|^{p}\mathrm{d}x\right)
-p\left(\|u\|_{E}^{2}
-\int_{\mathbb{R}^{N}}|u|^{p}\mathrm{d}x\right)\\
=&(2-p)\|u\|_{E}^{2}
<0.
\end{aligned}
\end{equation*}
Thus, $\Phi'(u)\not=0$ for $u\in \mathcal{N}$.

If $u\in \mathcal{N}$ and $I(u)=\bar{c}$,
from $\bar{c}$ is the minimum of $I$ on $\mathcal{N}$,
and then using the Lagrange multiplier theorem, we know that there exists $\lambda\in \mathbb{R}$ such that $I'(u)=\lambda \Phi'(u)$. So
\begin{equation*}
\begin{aligned}
\langle \lambda \Phi'(u),u\rangle =\langle I'(u),u\rangle=\Phi(u)=0.
\end{aligned}
\end{equation*}
This shows $\lambda=0$ and $I'(u)=0$. Thus, $u$ is a ground state solution for equation \eqref{P}.
\end{proof}

\end{document}
